\documentclass[11pt,a4paper]{article}
\usepackage[margin=1.8cm,headheight=48pt,headsep=12pt,footskip=22pt]{geometry}
\usepackage{fontspec}
\setmainfont{Noto Serif}
\setsansfont{Noto Sans}
\setmonofont{DejaVu Sans Mono}
\usepackage{amsmath,amssymb,mathtools}
\usepackage{enumitem,tabularx,booktabs,array,multicol}
\usepackage{xcolor,graphicx,tikz,fancyhdr,lastpage}
\usepackage{listings,tcolorbox,hyperref}
\hypersetup{colorlinks=true,linkcolor=blue!60!black,urlcolor=blue!60!black}
\definecolor{hcmusblue}{RGB}{28,55,111}
\definecolor{lightblue}{RGB}{237,242,250}
\definecolor{codebg}{RGB}{247,248,250}
\setlength{\parindent}{0pt}
\setlength{\parskip}{4pt}
\setlist[enumerate]{leftmargin=*,itemsep=5pt,topsep=4pt}
\setlist[itemize]{leftmargin=*,itemsep=2pt,topsep=2pt}
\lstset{language=C++,basicstyle=\ttfamily\footnotesize,backgroundcolor=\color{codebg},frame=single,rulecolor=\color{black!15},breaklines=true,showstringspaces=false,columns=fullflexible,keywordstyle=\color{hcmusblue}\bfseries,commentstyle=\color{green!40!black},numbers=left,numberstyle=\tiny\color{black!45},xleftmargin=1.5em,framexleftmargin=1.2em}
\pagestyle{fancy}
\fancyhf{}
\fancyhead[L]{\includegraphics[height=1.05cm,keepaspectratio]{assets/images/institution_logo.png}\hspace{0.45em}\begin{minipage}[c]{0.64\textwidth}\textbf{TRƯỜNG ĐẠI HỌC KHOA HỌC TỰ NHIÊN}\\[-1pt]\small ĐHQG TP. HỒ CHÍ MINH\end{minipage}}
\fancyhead[R]{\small\textcolor{hcmusblue}{TÀI LIỆU LUYỆN TẬP}}
\fancyfoot[L]{\scriptsize Biên soạn độc lập - Không phải đề thi chính thức}
\fancyfoot[R]{\scriptsize Trang \thepage}
\renewcommand{\headrulewidth}{0.5pt}
\renewcommand{\footrulewidth}{0.3pt}
\newcommand{\examtitle}[4]{\begin{center}{\Large\bfseries\color{hcmusblue}#1}\par\vspace{3pt}{\large\bfseries #2}\par\vspace{3pt}\small #3\quad\textbullet\quad #4\end{center}\vspace{5pt}\hrule\vspace{8pt}}
\newcommand{\studentline}{\begin{tabularx}{\textwidth}{X X}\textbf{Họ và tên:} \hrulefill & \textbf{MSSV:} \hrulefill\\[7pt]\textbf{Lớp:} \hrulefill & \textbf{Ngày làm:} \hrulefill\end{tabularx}\vspace{8pt}}
\newcommand{\points}[1]{\hfill\textbf{(#1 điểm)}}
\newtcolorbox{notebox}{colback=lightblue,colframe=hcmusblue!55,boxrule=0.5pt,arc=1.5mm,left=2mm,right=2mm,top=1mm,bottom=1mm}
\begin{document}
\examtitle{ĐỀ LUYỆN TẬP 03 - CƠ SỞ LẬP TRÌNH}{CƠ SỞ LẬP TRÌNH - HCMUS (theo lịch học tuần 1--5)}{90 phút - 10 điểm - C++11 - Không sử dụng tài liệu}{Giữa kỳ 1 - HK1 2026--2027}

\studentline
\begin{notebox}\textbf{Phạm vi:} Vòng lặp lồng nhau, hàm kiểm tra số nguyên tố, đệ quy cơ bản, mảng 1D và ma trận chuyển vị.\end{notebox}
\begin{enumerate}
\item Không chạy chương trình, hãy tính giá trị cuối cùng của \texttt{s} và giải thích số lần thực hiện câu lệnh cộng trong đoạn mã:\
\texttt{int s=0;}\\
\texttt{for (int i=1; i<=3; ++i)}\\
\texttt{\quad for (int j=1; j<=i; ++j) s += i+j;} \points{2}
\item Viết hàm \texttt{isPrime(n)} trả về \texttt{true} nếu $n$ là số nguyên tố, ngược lại trả về \texttt{false}. Sử dụng vòng lặp \texttt{for}; không cần dùng thư viện kiểm tra số nguyên tố. \points{2}
\item Viết hàm đệ quy \texttt{digitSum(n)} tính tổng chữ số của số nguyên không âm $n$. Nêu điều kiện dừng và minh họa với $n=5072$. \points{2}
\item Viết hàm tìm độ dài đoạn con liên tiếp tăng nghiêm ngặt dài nhất trong mảng. Đồng thời trả về vị trí bắt đầu của đoạn. Nếu có nhiều đoạn cùng độ dài, chọn đoạn xuất hiện trước. Với $n=0$, trả về độ dài 0 và \texttt{start=-1}. \points{2}
Prototype: \texttt{int longestIncreasingRun(const int a[], int n, int \&start);}
\item Viết hàm chuyển vị ma trận $R\times C$ sang ma trận $C\times R$. Kết quả đặt vào ma trận $T$; không yêu cầu chuyển vị tại chỗ. \points{2}
Prototype: \texttt{void transpose(const int A[][MAX], int T[][MAX], int R, int C);}
\end{enumerate}
\vfill\begin{center}\textit{--- Hết ---}\end{center}

\end{document}
